% GENERATED FILE. Edit canonical lessons, then run npm run generate:books.
% canonical-source-hash: fnv1a64:1a6c7d3edbe5f8a5
% canonical-lessons: PA-C67-nashta, PA-C67-anda, PA-C67-ghar, PA-C67-khirki, PA-C67-kursi

\chapter{Breakfast, Egg, House, Window, Chair}
\label{ch:pa-breakfast-egg-house-window-chair}
\hlchaptermodality{67}

\begin{chapteropening}
\textbf{By the end of this chapter:} \emph{I can say breakfast, an egg, a house, home, a window and a chair.}

Complete the last lesson of chapter 67: I can say breakfast, an egg, a house, home, a window and a chair.
\end{chapteropening}

\section[\texorpdfstring{\pa{ਨਾਸ਼ਤਾ} (nāshtā)}{ਨਾਸ਼ਤਾ (nāshtā)}]{\pa{ਨਾਸ਼ਤਾ} (nāshtā) — breakfast}

\label{lesson:PA-C67-nashta}



\begin{quote}
Before the new one: say the Punjabi for an orange, then the Punjabi for a lemon.
\end{quote}

\subsection*{You'll want to know: \pa{ਨਾਸ਼ਤਾ}}
\textbf{\pa{ਨਾਸ਼ਤਾ}} — \emph{nāshtā} — \textquotedblleft{}breakfast\textquotedblright{}.

\begin{cousinweb}
\textbf{\pa{ਨਾਸ਼ਤਾ}} (\emph{nāshtā}) is a loan, usually traced to Persian \emph{nāshtā}, \textquotedblleft{}not having eaten\textquotedblright{}. The meal is named for the fast it ends.
\end{cousinweb}

\begin{grammarlens}[title={What you've built}]
One more word for this chapter.
\end{grammarlens}

\subsection*{Your turn}
\begin{itemize}
\raggedright
  \item \emph{Say it:} \emph{nāshtā}
  \item \emph{Say it:} \emph{nāshtā}, once more
  \item \emph{Say it:} say \emph{nāshtā}
  \item \emph{Recall:} say the Punjabi for an orange, then the Punjabi for a lemon, then say \emph{nāshtā} again
\end{itemize}

\begin{tcolorbox}[breakable,colback=teal!4,colframe=teal!35!black,title={Before you move on}]
What does \pa{ਨਾਸ਼ਤਾ} mean? (\textquotedblleft{}Breakfast\textquotedblright{}.) Say it once more.
\end{tcolorbox}
\section[\texorpdfstring{\pa{ਆਂਡਾ} (āṇḍā)}{ਆਂਡਾ (āṇḍā)}]{\pa{ਆਂਡਾ} (āṇḍā) — an egg}

\label{lesson:PA-C67-anda}



\begin{quote}
Before the new one: say the Punjabi for a lemon, then the Punjabi for breakfast.
\end{quote}

\subsection*{You'll want to know: \pa{ਆਂਡਾ}}
\textbf{\pa{ਆਂਡਾ}} — \emph{āṇḍā} — \textquotedblleft{}an egg\textquotedblright{}.

\begin{cousinweb}
\textbf{\pa{ਆਂਡਾ}} (\emph{āṇḍā}) is a tadbhava, a word inherited from Sanskrit and reshaped by sound change. It comes from \emph{āṇḍa}, \textquotedblleft{}egg\textquotedblright{} (CDIAL 1111). CDIAL reports Mayrhofer's view, after Kuiper, that the word is of Munda origin.
\end{cousinweb}

\begin{grammarlens}[title={What you've built}]
One more word for this chapter.
\end{grammarlens}

\subsection*{Your turn}
\begin{itemize}
\raggedright
  \item \emph{Say it:} \emph{āṇḍā}
  \item \emph{Say it:} \emph{āṇḍā}, once more
  \item \emph{Say it:} say \emph{āṇḍā}
  \item \emph{Recall:} say the Punjabi for a lemon, then the Punjabi for breakfast, then say \emph{āṇḍā} again
\end{itemize}

\begin{tcolorbox}[breakable,colback=teal!4,colframe=teal!35!black,title={Before you move on}]
What does \pa{ਆਂਡਾ} mean? (\textquotedblleft{}An egg\textquotedblright{}.) Say it once more.
\end{tcolorbox}
\section[\texorpdfstring{\pa{ਘਰ} (ghar)}{ਘਰ (ghar)}]{\pa{ਘਰ} (ghar) — a house, home}

\label{lesson:PA-C67-ghar}



\begin{quote}
Before the new one: say the Punjabi for breakfast, then the Punjabi for an egg.
\end{quote}

\subsection*{You'll want to know: \pa{ਘਰ}}
\textbf{\pa{ਘਰ}} — \emph{ghar} — \textquotedblleft{}a house, home\textquotedblright{}.

\begin{cousinweb}
\textbf{\pa{ਘਰ}} (\emph{ghar}) is inherited. CDIAL lists it under \emph{ghara}, \textquotedblleft{}house\textquotedblright{}, and reports that it is generally taken as a popular form of Sanskrit \emph{gṛha}, \textquotedblleft{}house\textquotedblright{}, a link it calls phonetically difficult (CDIAL 4428).
\end{cousinweb}

\begin{grammarlens}[title={What you've built}]
One more word for this chapter.
\end{grammarlens}

\subsection*{Your turn}
\begin{itemize}
\raggedright
  \item \emph{Say it:} \emph{ghar}
  \item \emph{Say it:} \emph{ghar}, once more
  \item \emph{Say it:} say \emph{ghar}
  \item \emph{Recall:} say the Punjabi for breakfast, then the Punjabi for an egg, then say \emph{ghar} again
\end{itemize}

\begin{tcolorbox}[breakable,colback=teal!4,colframe=teal!35!black,title={Before you move on}]
What does \pa{ਘਰ} mean? (\textquotedblleft{}A house, home\textquotedblright{}.) Say it once more.
\end{tcolorbox}
\section[\texorpdfstring{\pa{ਖਿੜਕੀ} (khiṛkī)}{ਖਿੜਕੀ (khiṛkī)}]{\pa{ਖਿੜਕੀ} (khiṛkī) — a window}

\label{lesson:PA-C67-khirki}



\begin{quote}
Before the new one: say the Punjabi for an egg, then the Punjabi for a house.
\end{quote}

\subsection*{You'll want to know: \pa{ਖਿੜਕੀ}}
\textbf{\pa{ਖਿੜਕੀ}} — \emph{khiṛkī} — \textquotedblleft{}a window\textquotedblright{}.

\begin{cousinweb}
\textbf{\pa{ਖਿੜਕੀ}} (\emph{khiṛkī}) is a tadbhava, a word inherited from Sanskrit and reshaped by sound change. It comes from \emph{khaṭakkikā}, \textquotedblleft{}side door\textquotedblright{} (CDIAL 3770). CDIAL lists Punjabi \emph{khiṛkī}, \textquotedblleft{}back door, window\textquotedblright{}.
\end{cousinweb}

\begin{grammarlens}[title={What you've built}]
One more word for this chapter.
\end{grammarlens}

\subsection*{Your turn}
\begin{itemize}
\raggedright
  \item \emph{Say it:} \emph{khiṛkī}
  \item \emph{Say it:} \emph{khiṛkī}, once more
  \item \emph{Say it:} say \emph{khiṛkī}
  \item \emph{Recall:} say the Punjabi for an egg, then the Punjabi for a house, then say \emph{khiṛkī} again
\end{itemize}

\begin{tcolorbox}[breakable,colback=teal!4,colframe=teal!35!black,title={Before you move on}]
What does \pa{ਖਿੜਕੀ} mean? (\textquotedblleft{}A window\textquotedblright{}.) Say it once more.
\end{tcolorbox}
\section[\texorpdfstring{\pa{ਕੁਰਸੀ} (kursī)}{ਕੁਰਸੀ (kursī)}]{\pa{ਕੁਰਸੀ} (kursī) — a chair}

\label{lesson:PA-C67-kursi}



\begin{quote}
Before the new one: say the Punjabi for a house, then the Punjabi for a window.
\end{quote}

\subsection*{You'll want to know: \pa{ਕੁਰਸੀ}}
\textbf{\pa{ਕੁਰਸੀ}} — \emph{kursī} — \textquotedblleft{}a chair\textquotedblright{}.

\begin{cousinweb}
\textbf{\pa{ਕੁਰਸੀ}} (\emph{kursī}) is a loan from Arabic \emph{kursī}, \textquotedblleft{}chair, throne\textquotedblright{}, which reached Punjabi through Persian. Its Arabic root is \emph{k-r-s} (Lane).
\end{cousinweb}

\begin{grammarlens}[title={What you've built}]
One more word for this chapter.
\end{grammarlens}

\subsection*{Your turn}
\begin{itemize}
\raggedright
  \item \emph{Say it:} \emph{kursī}
  \item \emph{Say it:} \emph{kursī}, once more
  \item \emph{Say it:} say \emph{kursī}
  \item \emph{Recall:} say the Punjabi for a house, then the Punjabi for a window, then say \emph{kursī} again
\end{itemize}

\begin{tcolorbox}[breakable,colback=teal!4,colframe=teal!35!black,title={Before you move on}]
What does \pa{ਕੁਰਸੀ} mean? (\textquotedblleft{}A chair\textquotedblright{}.) Say it once more.
\end{tcolorbox}
